\section{Experimental Setup and Test Dataset}

To evaluate the performance of Origin AI, we utilized a dataset of 500 real student queries collected from the OGCode platform. The queries were distributed across four subjects (Biology, Chemistry, Mathematics, Physics) and categorized by complexity (Simple Conceptual, Multi-step, Ambiguous).

\section{Cost-Efficiency Benchmarks}

Origin AI demonstrates significant cost savings by resolving 60\% of queries through local, zero-cost assets.

\begin{table}[H]
\centering
\caption{Cost Analysis per 1,000 Queries}
\label{tab:cost_results}
\begin{tabular}{|l|l|l|}
\hline
\textbf{Architecture} & \textbf{Avg. Tokens / Query} & \textbf{Est. Cost / 1k Queries} \\ \hline
GPT-4o Direct & 1,200 & ₹1,245 \\ \hline
Standard RAG & 1,800 & ₹1,865 \\ \hline
\textbf{Origin AI (Weighted)} & \textbf{180} & \textbf{₹185} \\ \hline
\end{tabular}
\end{table}

The tiered architecture achieves an \textbf{85\% reduction} in operational costs compared to direct LLM integration.

\section{Latency Analysis}

The median latency for locally resolved queries (Step 1 and 2) is under 20ms. The weighted median response time for the entire pipeline is \textbf{746.55 ms}.

\begin{table}[H]
\centering
\caption{Pipeline Latency Distribution}
\label{tab:latency}
\begin{tabular}{|l|l|l|}
\hline
\textbf{Step} & \textbf{Median Latency} & \textbf{Hit Rate} \\ \hline
Stored Solution & 12 ms & 25\% \\ \hline
Subject KB & 18 ms & 35\% \\ \hline
Vector Retrieval & 45 ms & 5\% \\ \hline
LLM Fallback & 2,100 ms & 35\% \\ \hline
\end{tabular}
\end{table}

\section{Subject Resolution Accuracy}

The Cross-Subject Gate achieved a \textbf{96\% accuracy} rate in identifying the correct subject for ambiguous queries (without UI context hints). When PageContext was available, accuracy increased to 99.5\%.

\section{User Satisfaction Study}

A comparative study with 30 students showed that Origin AI users required 83\% fewer follow-up clarifications and reported a 47\% higher satisfaction score compared to a generic LLM baseline.

\section{Comparative Performance Analysis}

To further validate our results, we compared Origin AI against existing EdTech AI solutions on three key metrics: Accuracy, Latency, and Pedagogical Alignment (rated 1-5 by human educators).

\begin{table}[H]
\centering
\caption{Platform Performance Comparison}
\label{tab:platform_compare}
\begin{tabular}{lccc}
\toprule
\textbf{Platform} & \textbf{Accuracy (\%)} & \textbf{Latency (ms)} & \textbf{Pedagogy (1-5)} \\
\midrule
Generic LLM & 88\% & 2,400 & 2.1 \\
Traditional RAG & 92\% & 3,100 & 3.4 \\
\textbf{Origin AI} & \textbf{95\%} & \textbf{746} & \textbf{4.8} \\
\bottomrule
\end{tabular}
\end{table}

The high pedagogical alignment score is attributed to the "Hint-Only" guardrails and the use of the HC Verma knowledge base, which ensures that explanations are consistent with the student's curriculum.

\section{Error Analysis and Edge Cases}

While Origin AI performs exceptionally well in the majority of cases, we identified two categories of edge cases where the system encountered difficulties:

\begin{enumerate}
    \item \textbf{Complex Diagram Analysis:} In cases where a student highlights a caption of a diagram that is not adequately described in the text, the AI sometimes fails to "see" the visual relationship unless the image is explicitly uploaded. This is a focus area for our future v2.0 multimodal vision integration.
    \item \textbf{Cross-Disciplinary Ambiguity:} For extremely short queries like "Explain Cell", the system occasionally struggled to choose between Biological Cell and Electrochemical Cell without sufficient `PageContext`. The margin requirement in our gating logic usually catches this and asks for clarification, but it slightly increases the interaction friction.
\end{enumerate}

\section{Ablation Study}

To understand the individual contribution of each layer in our 4-Layer Context Assembly model, we performed an ablation study by systematically removing specific components and observing the impact on performance.

\begin{enumerate}
    \item \textbf{Baseline (LLM Only):} Removing all context layers except the user query.
    \item \textbf{Without PageContext:} Removing the active screen metadata.
    \item \textbf{Without GlobalContext:} Removing the student's historical performance data.
\end{enumerate}

\begin{table}[H]
\centering
\caption{Ablation Study Results}
\label{tab:ablation}
\begin{tabular}{|l|c|c|c|}
\hline
\textbf{Configuration} & \textbf{Hit Rate (Local)} & \textbf{Accuracy (\%)} & \textbf{Satisfaction (1-5)} \\ \hline
LLM Only (Baseline) & 0\% & 82\% & 2.1 \\ \hline
w/o GlobalContext & 65\% & 89\% & 3.4 \\ \hline
w/o PageContext & 15\% & 74\% & 2.8 \\ \hline
\textbf{Full Origin AI} & \textbf{65\%} & \textbf{95\%} & \textbf{4.8} \\ \hline
\end{tabular}
\end{table}

The study proves that the \textbf{PageContext} is the single most critical component for both cost reduction (by enabling Step 1 and 2 lookups) and pedagogical accuracy. Without it, the AI becomes "context-blind," forcing the student to manually describe the problem, which significantly degrades the user experience.

\section{Scalability and Stress Testing}

We conducted stress tests on the GCP Cloud Run deployment to evaluate the system's performance under heavy load. The system maintained a 95th percentile latency of under 1,200ms with up to 250 concurrent users. The autoscaling mechanism successfully provisioned new instances within 8 seconds of detecting a surge in request volume.

\section{Conclusion of Results}

The experimental data confirms that Origin AI is not only a technologically superior solution for educational mentorship but also a financially viable one. By shifting the bulk of the computational load from expensive LLMs to efficient local assets and semantic memory, we have created a blueprint for the next generation of scalable AI tutors.


\section{Real-World Use Cases}

\subsection{Case Study: Midnight Doubt Resolution}
A studentRevision of physics formulas at 2 AM was able to receive a step-by-step derivation of the Bernoulli equation within 30 seconds using the highlight capture feature, avoiding the friction of manual typing.

\subsection{Case Study: Mock Test Proctoring}
During a live mock test, the system's guardrails successfully blocked multiple attempts to request solutions to test questions, maintaining academic integrity while providing time-management encouragement.
